\begin{papercard}{tricco2018}
\cf{Stages} \stagetag{R} (reporting items for a review that maps a heterogeneous evidence base), \stagetag{A} (eligibility criteria and information sources for a non-intervention source set), \stagetag{D} (``data charting'': what is extracted, by whom, in what form), \stagetag{S} (why a scoping approach). PRISMA-ScR (\emph{Annals of Internal Medicine} 2018, 169:467--473; the PDF is the accepted manuscript from White Rose, 27 pages): 20 essential plus 2 optional reporting items for scoping reviews, adapted from the 27-item PRISMA 2009 checklist, with the consensus process that produced them. It is a reporting guideline, not a method; the explanation and elaboration is Supplement~2 and is not in the PDF.

\cf{What a scoping review is for} The abstract defines scoping reviews as following ``a systematic approach to map evidence on a topic; identify main concepts, theories and sources; and determine where the gaps are'' \pc{tricco2018}{7}{a type of knowledge synthesis, follow a systematic approach to map}. Objectives named in the introduction: examine the ``extent (i.e., size), range (i.e., variety) and nature (i.e., characteristics) of the evidence'' on a topic, decide whether a systematic review is worth doing, summarise a body of knowledge heterogeneous in methods or discipline, or identify gaps to plan research \pc{tricco2018}{8}{extent (i.e., size), range (i.e., variety) and nature (i.e., characteristics) of the evidence on}. Systematic reviews answer narrowly defined questions (does this intervention improve specified outcomes?); scoping reviews answer broader ones (what is known about this concept?) \pc{tricco2018}{9}{whereas scoping reviews are useful for answering much broader}. Because of that difference some PRISMA items do not fit and others are missing; the extension is also meant to apply to evidence maps \pc{tricco2018}{9}{This extension is also intended to be applicable to}. Before PRISMA-ScR, conduct guidance existed (JBI 2015/2017, building on Arksey and O'Malley 2005 and Levac et al.\ 2010) ``However, a reporting guideline for scoping reviews currently does not exist'' \pc{tricco2018}{8}{However, a reporting guideline for scoping reviews}.

\cf{Development} Protocol listed on the EQUATOR and PRISMA websites; two leads, two coordinators and a 4-member advisory board who did not score; the panel was meant to be representative by geography and stakeholder type. Modified Delphi: round 1 online survey (Qualtrics) of the modifications to PRISMA --- a broader question and search, optional risk-of-bias assessment and a consultation exercise, qualitative analysis --- rated on a 7-point Likert scale \pc{tricco2018}{10}{using a 7-point Likert scale}. An \emph{a priori} consensus rule of 85\,\% rating 6 or 7 marked agreement; below that an item was ``discrepant'' \pc{tricco2018}{11}{An 85\% consensus rule was selected a priori to signify agreement amongst the expert}. Round 2 had two arms: an in-person meeting in Toronto on 29 November 2016 under Chatham House rule with a modified nominal group technique and live re-scoring (16 attendees), and an online discussion on Conceptboard with re-scoring in Qualtrics; round 3 was a full-panel survey of items that had failed in both arms after six working-group calls. A workshop with 30 participants at the Global Evidence Summit (Cape Town, September 2017) tested the checklist.

\cf{Results of the consensus process} 37 invited, 31 completed round 1 and 24 completed all three rounds \pc{tricco2018}{13}{A total of 37 individuals were invited to participate}. In round 1 only 9 of 27 items reached agreement in the in-person arm and 8 of 27 in the e-Delphi arm; of 18 discrepant items 11 were re-scored in person, all passing except information sources (83\,\%) \pc{tricco2018}{14}{Of the 18 discrepant items from round 1, 11 were re-scored}; in the online arm 7 of 14 participated in discussion and 11 re-scored 19 items, of which 5 passed. Round 3 failed for two items: risk of bias across studies was deleted, and data collection process was kept but reworded to drop Arksey and O'Malley's optional consultation exercise, treated as a knowledge-translation activity that could accompany any synthesis \pc{tricco2018}{15}{was subsequently removed from the checklist, though item 10 was retained but}.

\cf{What changed from PRISMA} Final checklist: 20 items plus two optional \pc{tricco2018}{16}{The final checklist, with 20 items plus two optional items, is presented in Table 1}. Five PRISMA items were deemed not relevant (13 summary measures; 15 risk of bias across studies; 16 and 23 additional analyses; 22 risk of bias across studies results) \pc{tricco2018}{16}{Five items from the original PRISMA were deemed not relevant}. Risk of bias in individual studies (12) and its results (19) became \emph{optional} critical-appraisal items, because scoping reviews may include documents, blogs, websites, interviews and opinions and ``are not conducted to examine the risk of bias of the included sources'' \pc{tricco2018}{16}{opinions) and are not conducted to examine the risk of bias of the included sources}. Wording of every item was modified \pc{tricco2018}{17}{five items were removed (as they were deemed not}.

\cf{The checklist (Table 1)} \emph{Title} 1 identify as a scoping review. \emph{Abstract} 2 structured summary (background, objectives, eligibility, sources, charting methods, results, conclusions) \pc{tricco2018}{23}{charting methods, results and conclusions that relate to the review question(s) and objective(s)}. \emph{Introduction} 3 rationale, including why the question suits a scoping approach \pc{tricco2018}{23}{Explain why the review question(s)/objective(s) lend themselves to a scoping review approach}; 4 objectives with key elements (population, concepts, context). \emph{Methods} 5 protocol and registration; 6 eligibility criteria with a rationale \pc{tricco2018}{23}{Specify the characteristics of the sources of evidence (e.g., years considered, language, publication status) used as criteria for eligibility, and provide a rationale}; 7 information sources with coverage and last-search date; 8 search strategy ``for at least one database'' \pc{tricco2018}{23}{Present the full electronic search strategy for at least one database, including any limits used, such that it could be repeated}; 9 selection of sources; 10 data charting process, e.g.\ piloted forms and whether charting was done independently or in duplicate \pc{tricco2018}{24}{whether data charting was done independently, in duplicate}; 11 data items, ``list and define all variables'' \pc{tricco2018}{24}{List and define all variables for which data were sought and any assumptions and simplifications made}; 12 critical appraisal, if done \pc{tricco2018}{24}{If done, provide a rationale for conducting a critical appraisal of included sources of evidence}; 14 synthesis of results. \emph{Results} 17 numbers screened, assessed and included with reasons for exclusion and ideally a flow diagram \pc{tricco2018}{24}{Give numbers of sources of evidence screened, assessed for eligibility, and included in the review, with reasons for exclusions at each stage, ideally using a flow diagram}; 18 characteristics of sources with citations \pc{tricco2018}{24}{For each source of evidence, present characteristics for which data were charted and provide the citations}; 19 appraisal results; 20 results of individual sources; 21 synthesis. \emph{Discussion} 24 summary of evidence, ``including an overview of concepts, themes, and types of evidence'' \pc{tricco2018}{25}{Summarize the main results (including an overview of concepts, themes, and types of evidence available)}; 25 limitations; 26 conclusions. \emph{Funding} 27 funding of the included sources \emph{and} of the review itself \pc{tricco2018}{25}{Describe sources of funding for the included sources of evidence, as well as sources of funding for the scoping review}.

\cf{Mini-glossary} \emph{Charting} is the term for data extraction (Arksey and O'Malley; Levac; JBI) \pc{tricco2018}{25}{Charting – The process of data extraction in a scoping review is referred to as}. \emph{Sources of evidence} is deliberately broader than ``studies'' (research, expert opinion, policy documents) and is distinct from \emph{information sources}, the places where evidence is compiled (databases, social media platforms, websites) \pc{tricco2018}{25}{A more inclusive/ heterogeneous term is used to account for the fact that different types of evidence}. ``Critical appraisal'' replaces ``risk of bias'' for items 12 and 19 \pc{tricco2018}{25}{This terminology is used for items 12 and 19, instead of}.

\cf{Authors' statements on limits and impact} The guideline is said to be consistent with the JBI guidance; adherence ``will be evaluated in the future'', which is ``critical to measure its impact'' \pc{tricco2018}{17}{and that adherence to our reporting guideline will be evaluated in the future}; dissemination plans include short videos, webinars and one-page tip sheets.

\cf{Limits (card)} The panel was small (24 completed all rounds) and North-American-weighted; the round-2 online arm had half the discussion participation of the in-person arm, and the item-level percentages were not tested for reliability; the 85\,\% threshold and the ``not applicable'' status of five items are consensus decisions, not empirical findings; the PDF is an accepted manuscript, so page numbers and wording may differ from the published article, and the explanation with examples is not included; item 8 still asks only for ``at least one database'', a weaker demand than PRISMA 2020 and PRISMA-S (all sources, exactly as run); no guidance on non-bibliographic searches such as app stores.

\cf{Use for this thesis} (1) The thesis is neither a systematic review nor a meta-analysis; its literature part is closer to a scoping review or evidence map, so PRISMA-ScR is the closer reporting reference for the review sections, to be used with PRISMA-S for search details and PRISMA 2020 for the flow diagram. (2) ``Sources of evidence'' versus ``information sources'' gives vocabulary for the stage-A write-up: the apps and policies are \emph{sources of evidence}, the app stores and the policy web pages are \emph{information sources}. (3) Items 10--11 map directly to stage C and D: a piloted charting form, independent or duplicate charting, and listing and defining every variable are exactly the codebook, pilot and double-coding protocol; the thesis should state which of these were done. (4) Items 12 and 19 justify omitting risk-of-bias scoring while still stating that appraisal was not done. (5) Item 27 (funding of included sources) is a prompt to note commercial interests of app providers when analysing their policies. (6) Do not claim compliance; list adapted items, as with PRISMA 2020.
\end{papercard}
